\begin{helper}
Let \(U\) be the finite terminal coordinate set for a fixed history cell.
Let \(\mathcal T\) be the branch-label type and let \(\mathcal P\) be the
transcript-label type.  Let
\[
\mathsf{Realized}(\beta,\tau)
\]
mean that \((\beta,\tau)\) is a nonempty fixed support-pair cell, and let
\[
\mathsf{Comp}(A,\beta,\tau)
\]
mean that the complete terminal assignment \(A\subseteq U\) is compatible
with that realized cell.

Assume the branch choice and prefix scan are functional from a complete
terminal assignment: for every \(A\subseteq U\), if
\(\mathsf{Realized}(\beta,\tau)\) and
\(\mathsf{Comp}(A,\beta,\tau)\), then \(\beta\) is the branch label selected
from \(A\)'s complete terminal blue truth table and \(\tau\) is the
transcript returned by the deterministic prefix scan on \(A\).
Then a complete terminal assignment is compatible with at most one realized
branch/transcript cell.  Equivalently, if \(A\subseteq U\) is compatible
with two realized cells \((\beta,\tau)\) and \((\beta',\tau')\), then
\[
\beta=\beta'\qquad\text{and}\qquad \tau=\tau'.
\]
\end{helper}
\begin{proof}
Fix a complete terminal assignment \(A\subseteq U\) and suppose that it is
compatible with two realized cells \((\beta,\tau)\) and
\((\beta',\tau')\).  By the functionality hypothesis applied to the first
cell, the branch label \(\beta\) is exactly the branch selected from the
complete terminal blue truth table of \(A\), and the deterministic prefix
scan on \(A\) returns \(\tau\).  Applying the same functionality hypothesis
to the second cell shows that \(\beta'\) is the same selected branch label
from the same terminal assignment \(A\), and that the same scan on \(A\)
returns \(\tau'\).

The selected branch label is a function of \(A\), so the two branch labels
are equal:
\[
\beta=\beta'.
\]
The scan output is also a function of \(A\).  Since the same scan output is
simultaneously \(\tau\) and \(\tau'\), the transcript labels are equal:
\[
\tau=\tau'.
\]
Thus no complete terminal assignment can support two different realized
branch/transcript cells.  This is the precise disjoint-cylinder and
cell-to-residual-leaf injectivity fact needed by the red residual product
tree.
\end{proof}
